\documentclass[11pt]{ctexart}
\usepackage[a4paper,margin=1in]{geometry}
\usepackage{longtable,booktabs}
\title{Geant4-Agent Casebank Regression Report (EN)}
\author{Automated Pipeline}
\date{2026-03-08\_124606}
\begin{document}
\maketitle
\section{Unified Pass/Fail Table}
\begin{longtable}{p{0.06\linewidth}p{0.08\linewidth}p{0.07\linewidth}p{0.08\linewidth}p{0.07\linewidth}p{0.07\linewidth}p{0.07\linewidth}p{0.07\linewidth}p{0.09\linewidth}p{0.06\linewidth}p{0.21\linewidth}}
\toprule
ID & Domain & Style & Completeness & ExpInit & ActInit & ExpFinal & ActFinal & Missing(I->F) & Result & Fail Reason \\
\midrule
MTM01 & medical & fuzzy & missing & FAIL & FAIL & PASS & PASS & 12->0 & PASS & - \\
MTM02 & aerospace & precise & missing & FAIL & FAIL & PASS & PASS & 2->0 & PASS & - \\
MTM03 & industrial\_ndt & precise & missing & FAIL & FAIL & PASS & PASS & 12->0 & PASS & - \\
MTM04 & physics & precise & missing & FAIL & FAIL & PASS & PASS & 8->0 & PASS & - \\
MTM05 & nuclear\_engineering & fuzzy & missing & FAIL & FAIL & PASS & PASS & 5->0 & PASS & - \\
MTM06 & security\_screening & fuzzy & missing & FAIL & FAIL & PASS & PASS & 11->0 & PASS & - \\
MTM07 & semiconductor & precise & missing & FAIL & FAIL & PASS & PASS & 3->0 & PASS & - \\
MTM08 & education & fuzzy & missing & FAIL & FAIL & PASS & PASS & 11->0 & PASS & - \\
MTM09 & environmental\_radiation & fuzzy & missing & FAIL & FAIL & PASS & PASS & 3->0 & PASS & - \\
MTM10 & high\_energy\_physics & precise & missing & FAIL & FAIL & PASS & PASS & 10->0 & PASS & - \\
MTM11 & materials\_update & precise & full & PASS & PASS & PASS & PASS & 0->0 & PASS & - \\
MTM12 & overwrite\_reject & precise & full & PASS & PASS & PASS & PASS & 0->0 & PASS & - \\
MTM13 & medical\_detector & fuzzy & missing & FAIL & FAIL & PASS & PASS & 13->0 & PASS & - \\
MTM14 & industrial\_panel & precise & missing & FAIL & FAIL & PASS & PASS & 15->0 & PASS & - \\
MTM15 & shielding\_stack & precise & missing & FAIL & FAIL & PASS & PASS & 14->0 & PASS & - \\
MTM16 & nested\_target & precise & missing & FAIL & FAIL & PASS & PASS & 13->0 & PASS & - \\
MTM17 & coaxial\_shield & fuzzy & missing & FAIL & FAIL & PASS & PASS & 10->0 & PASS & - \\
MTM18 & boolean\_cavity & precise & missing & FAIL & FAIL & PASS & PASS & 17->0 & PASS & - \\
MTM19 & ring\_complete & precise & full & PASS & PASS & PASS & PASS & 0->0 & PASS & - \\
MTM20 & grid\_complete & precise & full & PASS & PASS & PASS & PASS & 0->0 & PASS & - \\
MTM21 & layered\_detector & fuzzy & missing & FAIL & FAIL & PASS & PASS & 14->0 & PASS & - \\
MTM22 & embedded\_sensor & precise & full & PASS & PASS & PASS & PASS & 0->0 & PASS & - \\
MTM23 & layered\_shield\_complete & precise & full & PASS & PASS & PASS & PASS & 0->0 & PASS & - \\
MTM24 & boolean\_complete & precise & full & PASS & PASS & PASS & PASS & 0->0 & PASS & - \\
\bottomrule
\end{longtable}
\section{Baseline Comparison}
\begin{itemize}
\item pass\_rate=1.000
\item initial\_alignment\_rate=1.000
\item final\_alignment\_rate=1.000
\item avg\_turns\_per\_case=2.458
\item missing\_reduction\_rate=1.000
\item confirm\_apply\_success\_rate=1.000
\item rollback\_after\_confirm\_rate=0.000
\item internal\_leak\_turn\_rate=0.068
\item repeated\_assistant\_case\_rate=0.000
\item stale\_confirm\_case\_rate=0.000
\item delta\_pass\_rate=+0.050
\item regressions=none
\item improvements=none
\end{itemize}
\appendix
\section{Appendix: Full Dialogues and Naturalized Internal Chains}
\subsection{Case MTM01 (medical)}
\textbf{Dialogue:}
\begin{itemize}
\item User: 我先要一个用于医学成像的初步辐照场景，细节还没定，你先告诉我还缺什么。
\item Assistant: 目前我们已经确定了使用盒体作为几何结构，接下来还需要补充一些具体信息：

- **几何**：需要盒体在x、y、z三个方向上的具体尺寸。
- **材料**：需要指定体积与材料的绑定关系。
- **源**：需要选择源的类型（例如粒子束、放射源等）。

请随时告诉我这些细节，我们可以一步步完善这个辐照场景。
\item User: 采用 300 mm x 300 mm x 300 mm 的水箱，材料 ?；源用束流 gamma，能量 6 MeV，位置 (0,0,-400) mm，方向 +z；物理列表 QBBC；输出 root。
\item Assistant: 我检测到这会覆盖已经确认的值。待修改内容：几何结构类型：unknown -> single\_box；盒体 x 方向尺寸：None -> 300.0。如果要应用这些修改，请直接回复“确认”；如果要保留原值，请直接说“保持原值”。
\item User: 确认
\item Assistant: 配置已完成。当前已确认的几何参数包括：盒体在 x、y、z 方向上的尺寸以及几何结构类型。
\end{itemize}
\textbf{Internal Chain (Natural Language):}
\begin{itemize}
\item Turn 1: action=summarize\_progress, phase=geometry, complete=False.
\item Missing fields: geometry.params.module\_x, geometry.params.module\_y, geometry.params.module\_z, materials.volume\_material\_map, source.type, source.particle, source.energy, source.position, source.direction, physics.physics\_list, output.format, output.path.
\item LLM used=True; fallback=none; backend=llm\_semantic\_frame+runtime\_semantic.
\item Turn 2: action=confirm\_overwrite, phase=geometry, complete=False.
\item Missing fields: geometry.params.module\_x, geometry.params.module\_y, geometry.params.module\_z.
\item LLM used=True; fallback=none; backend=llm\_slot\_frame+runtime\_semantic.
\item Turn 3: action=finalize, phase=materials, complete=True.
\item Missing fields: none.
\item LLM used=True; fallback=none; backend=llm\_semantic\_frame+runtime\_semantic+confirmed\_pending\_overwrite.
\end{itemize}
\subsection{Case MTM02 (aerospace)}
\textbf{Dialogue:}
\begin{itemize}
\item User: I need a first-pass proton irradiation setup for a silicon device; keep it simple for now.
\item Assistant: Geometry: updated geometry structure, box size along x, box size along y. Materials: updated selected materials, volume-to-material binding. Source: updated source type, particle type, source energy. Physics: updated physics list. Output: still needs output format, output path.
\item User: Use a silicon box 50 mm x 50 mm x 5 mm, beam source proton 100 MeV from (0,0,-100) mm along +z, physics FTFP\_BERT, output json.
\item Assistant: You're about to overwrite a previously confirmed value.

Pending change:
- Box size along x: 1000.0 → 50.0
- Box size along y: 1000.0 → 50.0

Reply **'confirm'** to apply these changes, or **'keep existing'** to keep the current values.
\item User: confirm
\item Assistant: Configuration complete. Confirmed items: source type, box size along z, box size along y, and box size along x.
\end{itemize}
\textbf{Internal Chain (Natural Language):}
\begin{itemize}
\item Turn 1: action=summarize\_progress, phase=materials, complete=False.
\item Missing fields: output.format, output.path.
\item LLM used=True; fallback=none; backend=llm\_slot\_frame+runtime\_semantic.
\item Turn 2: action=confirm\_overwrite, phase=source, complete=False.
\item Missing fields: none.
\item LLM used=True; fallback=none; backend=llm\_slot\_frame+runtime\_semantic.
\item Turn 3: action=finalize, phase=physics, complete=True.
\item Missing fields: none.
\item LLM used=True; fallback=none; backend=llm\_semantic\_frame+runtime\_semantic+confirmed\_pending\_overwrite.
\end{itemize}
\subsection{Case MTM03 (industrial\_ndt)}
\textbf{Dialogue:}
\begin{itemize}
\item User: 我要做一个 gamma 穿透铁靶的小测试，但我还没把完整参数说完。
\item Assistant: 目前我们已确定使用铁作为材料，并选择 gamma 粒子作为源类型。

接下来还需要你补充几何部分的信息，包括结构类型、盒体在 x、y、z 方向的尺寸，以及将体积与材料绑定的设置。
\item User: 请设为半径 40 mm、半长 80 mm 的铁圆柱，材料 ?；点源 gamma，能量 2 MeV，位置 (0,0,-200) mm，方向 +z；物理列表 Shielding；输出 csv。
\item Assistant: 配置已完成！当前已确认的设置包括：体积与材料绑定、输出路径、输出格式和物理列表。
\end{itemize}
\textbf{Internal Chain (Natural Language):}
\begin{itemize}
\item Turn 1: action=summarize\_progress, phase=geometry, complete=False.
\item Missing fields: geometry.structure, geometry.params.module\_x, geometry.params.module\_y, geometry.params.module\_z, materials.volume\_material\_map, source.type, source.energy, source.position, source.direction, physics.physics\_list, output.format, output.path.
\item LLM used=True; fallback=none; backend=llm\_slot\_frame+runtime\_semantic.
\item Turn 2: action=finalize, phase=materials, complete=True.
\item Missing fields: none.
\item LLM used=True; fallback=none; backend=llm\_slot\_frame+runtime\_semantic.
\end{itemize}
\subsection{Case MTM04 (physics)}
\textbf{Dialogue:}
\begin{itemize}
\item User: Gamma point source 1 MeV at (0,0,-100) toward +z; please build the rest.
\item Assistant: Geometry: still needs geometry structure, box size along x. Materials: still needs volume-to-material binding. Source: updated source type, particle type, source energy. Physics: still needs physics list. Output: still needs output format, output path.
\item User: Set copper box 1 m x 1 m x 1 m, physics FTFP\_BERT, output hdf5.
\item Assistant: Your configuration is complete. The following items have been confirmed: volume-to-material binding, output path, output format, and physics list.
\end{itemize}
\textbf{Internal Chain (Natural Language):}
\begin{itemize}
\item Turn 1: action=summarize\_progress, phase=geometry, complete=False.
\item Missing fields: geometry.structure, geometry.params.module\_x, geometry.params.module\_y, geometry.params.module\_z, materials.volume\_material\_map, physics.physics\_list, output.format, output.path.
\item LLM used=True; fallback=none; backend=llm\_slot\_frame+runtime\_semantic.
\item Turn 2: action=finalize, phase=materials, complete=True.
\item Missing fields: none.
\item LLM used=True; fallback=none; backend=llm\_slot\_frame+runtime\_semantic.
\end{itemize}
\subsection{Case MTM05 (nuclear\_engineering)}
\textbf{Dialogue:}
\begin{itemize}
\item User: 我想做一个中子屏蔽的概念验证，现在先给一个缺参开头。
\item Assistant: 目前我们已经确定了几个关键部分：几何结构、材料、粒子源类型、物理列表以及输出设置。

接下来还需要补充一些具体参数：盒体在x、y、z三个方向上的尺寸，以及源的发射能量和位置。

请随时告诉我这些信息，我们可以继续推进配置。
\item User: 使用 200 mm x 200 mm x 200 mm 的铅块，材料 ?；点源 neutron，能量 14 MeV，位置 (0,0,-50) mm，方向 +z；物理列表 QGSP\_BERT\_HP；输出 xml。
\item Assistant: 我检测到这会覆盖已经确认的值。待修改内容：几何结构类型：unknown -> single\_box；已选材料：['G4\_Cu'] -> ['G4\_Pb']。如果要应用这些修改，请直接回复“确认”；如果要保留原值，请直接说“保持原值”。
\item User: 确认
\item Assistant: 配置已完成。当前已确认的参数包括：体积与材料绑定、输出路径、盒体 z 方向尺寸、盒体 y 方向尺寸。
\end{itemize}
\textbf{Internal Chain (Natural Language):}
\begin{itemize}
\item Turn 1: action=summarize\_progress, phase=geometry, complete=False.
\item Missing fields: geometry.params.module\_x, geometry.params.module\_y, geometry.params.module\_z, source.energy, source.position.
\item LLM used=True; fallback=none; backend=llm\_slot\_frame+runtime\_semantic.
\item Turn 2: action=confirm\_overwrite, phase=geometry, complete=False.
\item Missing fields: geometry.params.module\_x, geometry.params.module\_y, geometry.params.module\_z.
\item LLM used=True; fallback=none; backend=llm\_slot\_frame+runtime\_semantic.
\item Turn 3: action=finalize, phase=materials, complete=True.
\item Missing fields: none.
\item LLM used=True; fallback=none; backend=llm\_slot\_frame+runtime\_semantic+confirmed\_pending\_overwrite.
\end{itemize}
\subsection{Case MTM06 (security\_screening)}
\textbf{Dialogue:}
\begin{itemize}
\item User: Need a quick cargo-screening toy setup with forward photons, but I have not fixed the target shape, material, source energy, or output format yet.
\item Assistant: Geometry: updated geometry structure; still needs box size along x, box size along y. Materials: still needs volume-to-material binding. Source: updated particle type; still needs source type, source energy. Physics: still needs physics list. Output: still needs output format, output path.
\item User: Use steel cylinder target radius 60 mm half-length 120 mm, beam gamma 3 MeV from (0,0,-500) mm to +z, physics Shielding, output root.
\item Assistant: This would overwrite an already confirmed value. Pending change: geometry structure: unknown -> single\_tubs; box size along x: None -> 1000.0. Reply 'confirm' to apply it, or say 'keep existing' to keep the current value.
\item User: confirm
\item Assistant: Configuration complete. Confirmed items: volume-to-material binding, half length, radius, and box size along z.
\end{itemize}
\textbf{Internal Chain (Natural Language):}
\begin{itemize}
\item Turn 1: action=summarize\_progress, phase=geometry, complete=False.
\item Missing fields: geometry.params.module\_x, geometry.params.module\_y, geometry.params.module\_z, materials.volume\_material\_map, source.type, source.energy, source.position, source.direction, physics.physics\_list, output.format, output.path.
\item LLM used=True; fallback=none; backend=llm\_slot\_frame+runtime\_semantic.
\item Turn 2: action=confirm\_overwrite, phase=geometry, complete=False.
\item Missing fields: geometry.params.module\_x, geometry.params.module\_y, geometry.params.module\_z.
\item LLM used=True; fallback=none; backend=llm\_slot\_frame+runtime\_semantic.
\item Turn 3: action=finalize, phase=materials, complete=True.
\item Missing fields: none.
\item LLM used=True; fallback=none; backend=llm\_semantic\_frame+runtime\_semantic+confirmed\_pending\_overwrite.
\end{itemize}
\subsection{Case MTM07 (semiconductor)}
\textbf{Dialogue:}
\begin{itemize}
\item User: 我想搭一个硅薄片的中子照射测试，但先只给一个简述。
\item Assistant: 目前我们已经完成了几何结构（盒体尺寸）、材料（硅）和粒子源（中子能量与位置）的设置。

接下来还需要你补充以下信息：
- 物理过程列表
- 输出文件格式
- 输出文件保存路径

请告诉我这些内容，我们就可以继续了。
\item User: 使用 20 mm x 20 mm x 1 mm 的硅片，材料 ?；各向同性 neutron，能量 2 MeV，位置 (0,0,0) mm；物理列表 QGSP\_BERT\_HP；输出 json。
\item Assistant: 我注意到这次修改会覆盖已经确认的参数。具体变动如下：

- 盒体 x 方向尺寸：1000.0 → 20.0
- 盒体 y 方向尺寸：1000.0 → 20.0

请确认你的选择：
• 如果要应用这些修改，请直接回复“确认”。
• 如果要保留原来的数值，请直接回复“保持原值”。
\item User: 确认
\item Assistant: 配置已完成。当前已确认的参数包括：源位置、源能量、源类型以及盒体 z 方向尺寸。
\end{itemize}
\textbf{Internal Chain (Natural Language):}
\begin{itemize}
\item Turn 1: action=summarize\_progress, phase=materials, complete=False.
\item Missing fields: physics.physics\_list, output.format, output.path.
\item LLM used=True; fallback=none; backend=llm\_slot\_frame+runtime\_semantic.
\item Turn 2: action=confirm\_overwrite, phase=source, complete=False.
\item Missing fields: none.
\item LLM used=True; fallback=none; backend=llm\_slot\_frame+runtime\_semantic.
\item Turn 3: action=finalize, phase=physics, complete=True.
\item Missing fields: none.
\item LLM used=True; fallback=none; backend=llm\_semantic\_frame+runtime\_semantic+confirmed\_pending\_overwrite.
\end{itemize}
\subsection{Case MTM08 (education)}
\textbf{Dialogue:}
\begin{itemize}
\item User: For classroom demo, I want gamma attenuation in copper; you can ask only essential missing parts.
\item Assistant: Geometry: still needs geometry structure, box size along x. Materials: updated selected materials; still needs volume-to-material binding. Source: updated particle type; still needs source type, source energy. Physics: updated physics list. Output: still needs output format, output path.
\item User: Okay, copper box 100 mm x 100 mm x 100 mm, point gamma 1 MeV at (0,0,-200) mm, direction +z, physics FTFP\_BERT, output csv.
\item Assistant: You've already confirmed the physics list as "Shielding." Changing it to "FTFP\_BERT" would overwrite that.

Please reply **'confirm'** to apply the change to FTFP\_BERT, or **'keep existing'** to keep Shielding.
\item User: confirm
\item Assistant: Configuration complete. The following items are now confirmed: physics list, output path, output format, and source direction.
\end{itemize}
\textbf{Internal Chain (Natural Language):}
\begin{itemize}
\item Turn 1: action=summarize\_progress, phase=geometry, complete=False.
\item Missing fields: geometry.structure, geometry.params.module\_x, geometry.params.module\_y, geometry.params.module\_z, materials.volume\_material\_map, source.type, source.energy, source.position, source.direction, output.format, output.path.
\item LLM used=True; fallback=none; backend=llm\_slot\_frame+runtime\_semantic.
\item Turn 2: action=confirm\_overwrite, phase=materials, complete=False.
\item Missing fields: none.
\item LLM used=True; fallback=none; backend=llm\_slot\_frame+runtime\_semantic.
\item Turn 3: action=finalize, phase=source, complete=True.
\item Missing fields: none.
\item LLM used=True; fallback=none; backend=llm\_semantic\_frame+runtime\_semantic+confirmed\_pending\_overwrite.
\end{itemize}
\subsection{Case MTM09 (environmental\_radiation)}
\textbf{Dialogue:}
\begin{itemize}
\item User: 我想先搭一个空气体积中的低能 gamma 传播场景，细节稍后补。
\item Assistant: 目前已经完成了几何、材料和源的基本设置。接下来还需要你补充物理列表、输出格式和输出路径这三项信息。
\item User: 设为 500 mm x 500 mm x 500 mm 的空气盒，材料 ??；点源 gamma，能量 0.5 MeV，位置 (0,0,-300) mm，方向 +z；物理列表 QBBC；输出 root。
\item Assistant: 我注意到这次修改会覆盖你已经确认过的参数。具体变动如下：

- 盒体 x 方向尺寸：从 1000.0 改为 500.0
- 盒体 y 方向尺寸：从 1000.0 改为 500.0

请确认你的选择：
如果应用这些修改，请回复“确认”；
如果想保留原来的数值，请回复“保持原值”。
\item User: 确认
\item Assistant: 配置已完成。当前已确认的参数包括：源位置、源能量、盒体 z 方向尺寸、盒体 y 方向尺寸。
\end{itemize}
\textbf{Internal Chain (Natural Language):}
\begin{itemize}
\item Turn 1: action=summarize\_progress, phase=materials, complete=False.
\item Missing fields: physics.physics\_list, output.format, output.path.
\item LLM used=True; fallback=none; backend=llm\_slot\_frame+runtime\_semantic.
\item Turn 2: action=confirm\_overwrite, phase=source, complete=False.
\item Missing fields: none.
\item LLM used=True; fallback=none; backend=llm\_slot\_frame+runtime\_semantic.
\item Turn 3: action=finalize, phase=physics, complete=True.
\item Missing fields: none.
\item LLM used=True; fallback=none; backend=llm\_semantic\_frame+runtime\_semantic+confirmed\_pending\_overwrite.
\end{itemize}
\subsection{Case MTM10 (high\_energy\_physics)}
\textbf{Dialogue:}
\begin{itemize}
\item User: Set geometry as a sphere target; I have not decided radius or source details yet.
\item Assistant: Here's what we have so far:

**Geometry:** You've set the structure to a sphere. We still need the sphere's **radius**.

**Materials:** We need to assign a material to your volume.

**Source:** We need the source type, particle type, and particle energy.

**Physics \& Output:** We also need to select a physics list and set the output format and file path.
\item User: Sphere target radius 80 mm, lead material, electron point source 5 MeV at (0,0,-200) mm toward +z, physics FTFP\_BERT, output hdf5.
\item Assistant: Your configuration is complete. The following items are confirmed: volume-to-material binding, output path, output format, and physics list.
\end{itemize}
\textbf{Internal Chain (Natural Language):}
\begin{itemize}
\item Turn 1: action=summarize\_progress, phase=geometry, complete=False.
\item Missing fields: geometry.params.child\_rmax, materials.volume\_material\_map, source.type, source.particle, source.energy, source.position, source.direction, physics.physics\_list, output.format, output.path.
\item LLM used=True; fallback=none; backend=llm\_slot\_frame+runtime\_semantic.
\item Turn 2: action=finalize, phase=materials, complete=True.
\item Missing fields: none.
\item LLM used=True; fallback=none; backend=llm\_slot\_frame+runtime\_semantic.
\end{itemize}
\subsection{Case MTM11 (materials\_update)}
\textbf{Dialogue:}
\begin{itemize}
\item User: Create a complete setup: 100 mm x 100 mm x 100 mm copper box, point gamma 1 MeV at (0,0,-200) mm toward +z, physics QBBC, output json.
\item Assistant: Current committed config is: physics list FTFP\_BERT, output format json, source type point, particle gamma. If you want changes, please specify which field to update.
\item User: Change the target material to aluminum and change the output format to root.
\item Assistant: You're about to overwrite two confirmed settings:

- **Selected materials**: Copper → Aluminum
- **Output format**: JSON → ROOT

Please reply **"confirm"** to apply these changes, or **"keep existing"** to leave the current values as they are.
\item User: confirm
\item Assistant: Configuration complete. The following items are confirmed: your selected materials, output format, output path, and volume-to-material bindings.
\end{itemize}
\textbf{Internal Chain (Natural Language):}
\begin{itemize}
\item Turn 1: action=finalize, phase=materials, complete=True.
\item Missing fields: none.
\item LLM used=True; fallback=none; backend=llm\_slot\_frame+runtime\_semantic.
\item Turn 2: action=confirm\_overwrite, phase=source, complete=False.
\item Missing fields: none.
\item LLM used=True; fallback=none; backend=llm\_slot\_frame+runtime\_semantic.
\item Turn 3: action=finalize, phase=physics, complete=True.
\item Missing fields: none.
\item LLM used=True; fallback=none; backend=llm\_semantic\_frame+runtime\_semantic+confirmed\_pending\_overwrite.
\end{itemize}
\subsection{Case MTM12 (overwrite\_reject)}
\textbf{Dialogue:}
\begin{itemize}
\item User: 请创建一个完整配置：100 mm x 100 mm x 100 mm 的铜盒，材料 ?；点源 gamma，能量 1 MeV，位置 (0,0,-200) mm，方向 +z；物理列表 QBBC；输出 json。
\item Assistant: 配置已完成。当前已确认的设置包括：体积与材料绑定、输出路径和输出格式。
\item User: 把材料改成 ?，并把输出格式改成 root。
\item Assistant: 检测到您想要修改输出格式，这将会覆盖已确认的当前值（json 改为 root）。请确认您的选择：直接回复“确认”以应用修改，或回复“保持原值”以保留原来的设置。
\item User: 保持原值
\item Assistant: 好的，输出格式仍保持为 json，没有改动。
\end{itemize}
\textbf{Internal Chain (Natural Language):}
\begin{itemize}
\item Turn 1: action=finalize, phase=materials, complete=True.
\item Missing fields: none.
\item LLM used=True; fallback=none; backend=llm\_slot\_frame+runtime\_semantic.
\item Turn 2: action=confirm\_overwrite, phase=source, complete=False.
\item Missing fields: geometry.ask.boolean.solid\_a\_size, geometry.ask.boolean.solid\_b\_size.
\item LLM used=True; fallback=none; backend=llm\_slot\_frame+runtime\_semantic.
\item Turn 3: action=reject\_overwrite, phase=physics, complete=True.
\item Missing fields: none.
\item LLM used=True; fallback=none; backend=llm\_semantic\_frame+runtime\_semantic+rejected\_pending\_overwrite.
\end{itemize}
\subsection{Case MTM13 (medical\_detector)}
\textbf{Dialogue:}
\begin{itemize}
\item User: I want a ring of detector modules around a small source region, but I have not fixed the ring radius or module size yet.
\item Assistant: Your geometry structure is set. Next, we need to specify the materials for your volumes, define your particle source, and choose the physics and output settings.
\item User: Use a ring of 12 modules, each 8 mm x 10 mm x 2 mm, radius 40 mm, clearance 1 mm, material water; point gamma source 0.662 MeV at (0,0,0) mm; physics QBBC; output json.
\item Assistant: You are about to overwrite a confirmed geometry configuration.

**Current Geometry:**
*   Module: Box with dimensions 7.83 mm x 4.72 mm x 6.86 mm
*   Ring: 4 modules, radius 67.49 mm, clearance 0.19 mm

**New Geometry:**
*   Module: Box with dimensions 8.0 mm x 10.0 mm x 2.0 mm
*   Ring: 12 modules, radius 40.0 mm, clearance 1.0 mm

Please reply **'confirm'** to apply this change, or **'keep existing'** to keep the current values.
\item User: confirm
\item Assistant: Configuration complete. The following geometry items are confirmed: ring clearance, ring radius, ring module count, and ring module size.
\end{itemize}
\textbf{Internal Chain (Natural Language):}
\begin{itemize}
\item Turn 1: action=summarize\_progress, phase=materials, complete=False.
\item Missing fields: materials.volume\_material\_map, source.type, source.particle, source.energy, source.position, source.direction, physics.physics\_list, output.format, output.path, geometry.ask.ring.module\_size, geometry.ask.ring.count, geometry.ask.ring.radius, geometry.ask.ring.clearance.
\item LLM used=True; fallback=none; backend=llm\_semantic\_frame+runtime\_semantic.
\item Turn 2: action=confirm\_overwrite, phase=source, complete=False.
\item Missing fields: none.
\item LLM used=True; fallback=none; backend=llm\_slot\_frame+runtime\_semantic.
\item Turn 3: action=finalize, phase=physics, complete=True.
\item Missing fields: none.
\item LLM used=True; fallback=none; backend=llm\_semantic\_frame+runtime\_semantic+confirmed\_pending\_overwrite.
\end{itemize}
\subsection{Case MTM14 (industrial\_panel)}
\textbf{Dialogue:}
\begin{itemize}
\item User: 我想做一个二维阵列探测板，但 pitch 和材料还没定。
\item Assistant: 好的，目前配置进度如下：

**已完成：**
*   **几何**：已确定使用“二维阵列”结构。

**待补充：**
*   **几何**：阵列模块的尺寸、X方向的阵列数量。
*   **材料**：体积与材料的绑定关系。
*   **源**：源类型、粒子类型、能量和位置。
*   **物理**：物理过程列表。
*   **输出**：输出文件的格式和保存路径。

接下来，我们首先来确定阵列模块的具体尺寸吧。
\item User: 设成 4 x 3 的盒体阵列，每个模块 10 mm x 10 mm x 2 mm，pitch\_x 12 mm，pitch\_y 14 mm，间隙 1 mm，材料 硅；点源 gamma 0.511 MeV，位置 (0,0,-80) mm，方向 +z；物理列表 QBBC；输出 root。
\item Assistant: 我注意到这将会覆盖已确认的几何参数。具体修改如下：

*   模块尺寸从 (7.83, 4.72, 6.86) mm 改为 (10.0, 10.0, 2.0) mm。
*   X方向间距从 11.47 mm 改为 12.0 mm。
*   Y方向间距从 8.35 mm 改为 14.0 mm。
*   阵列间隙从 1.82 mm 改为 1.0 mm。

请确认：**回复“确认”将应用以上修改；回复“保持原值”则保留现有配置。**
\item User: 确认
\item Assistant: 配置已完成。当前已确认的参数包括：y方向阵列数量、x方向阵列数量、阵列间隙以及y方向的pitch。
\end{itemize}
\textbf{Internal Chain (Natural Language):}
\begin{itemize}
\item Turn 1: action=summarize\_progress, phase=materials, complete=False.
\item Missing fields: materials.volume\_material\_map, source.type, source.particle, source.energy, source.position, source.direction, physics.physics\_list, output.format, output.path, geometry.ask.grid.module\_size, geometry.ask.grid.count\_x, geometry.ask.grid.count\_y, geometry.ask.grid.pitch\_x, geometry.ask.grid.pitch\_y, geometry.ask.grid.clearance.
\item LLM used=True; fallback=none; backend=llm\_semantic\_frame+runtime\_semantic.
\item Turn 2: action=confirm\_overwrite, phase=source, complete=False.
\item Missing fields: geometry.ask.grid.count\_x, geometry.ask.grid.count\_y.
\item LLM used=True; fallback=none; backend=llm\_slot\_frame+runtime\_semantic.
\item Turn 3: action=finalize, phase=physics, complete=True.
\item Missing fields: none.
\item LLM used=True; fallback=none; backend=llm\_semantic\_frame+runtime\_semantic+confirmed\_pending\_overwrite.
\end{itemize}
\subsection{Case MTM15 (shielding\_stack)}
\textbf{Dialogue:}
\begin{itemize}
\item User: Build a stacked absorber sandwich along z; I will provide the layer sizes next.
\item Assistant: You've set up a stacked absorber sandwich geometry along the z-axis. Here's what's complete and what's still needed:

**Completed:**
*   **Geometry:** The overall structure is set.

**Still Needed:**
*   **Geometry:** The footprint and thicknesses for each stack layer.
*   **Materials:** Which material to assign to each volume.
*   **Source:** The source type, particle type, and its energy.
*   **Physics:** The physics list to use.
*   **Output:** The format and file path for your results.
\item User: Use stacked layers with footprint 20 mm x 20 mm, thicknesses 2 mm, 4 mm, 6 mm, layer clearance 0.5 mm, outer box 30 mm x 30 mm x 20 mm, container clearance 1 mm, material aluminum; beam gamma 1.25 MeV from (0,0,-120) mm along +z; physics Shielding; output csv.
\item Assistant: You are about to overwrite a confirmed geometry configuration.

**Current Geometry:**
*   **Stack:** Footprint \textasciitilde{}16.33 mm x 10.28 mm, Thicknesses [\textasciitilde{}6.68, \textasciitilde{}1.19, \textasciitilde{}5.59] mm, Clearance \textasciitilde{}0.73 mm.
*   **Container Box:** \textasciitilde{}18.07 mm x \textasciitilde{}45.52 mm x \textasciitilde{}16.66 mm, Clearance \textasciitilde{}0.87 mm.

**New Geometry:**
*   **Stack:** Footprint 20.0 mm x 20.0 mm, Thicknesses [2.0, 4.0, 6.0] mm, Clearance 0.5 mm.
*   **Container Box:** 30.0 mm x 30.0 mm x 20.0 mm, Clearance 1.0 mm.

Reply **'confirm'** to apply this change, or **'keep existing'** to keep the current values.
\item User: confirm
\item Assistant: Your configuration is complete. The following geometry items are confirmed: stack container clearance, stack container size, stack layer clearance, and stack layer thicknesses.
\end{itemize}
\textbf{Internal Chain (Natural Language):}
\begin{itemize}
\item Turn 1: action=summarize\_progress, phase=materials, complete=False.
\item Missing fields: materials.volume\_material\_map, source.type, source.particle, source.energy, source.position, source.direction, physics.physics\_list, output.format, output.path, geometry.ask.stack.footprint, geometry.ask.stack.thicknesses, geometry.ask.stack.layer\_clearance, geometry.ask.stack.parent\_size, geometry.ask.stack.nest\_clearance.
\item LLM used=True; fallback=none; backend=llm\_semantic\_frame+runtime\_semantic.
\item Turn 2: action=confirm\_overwrite, phase=source, complete=False.
\item Missing fields: none.
\item LLM used=True; fallback=none; backend=llm\_slot\_frame+runtime\_semantic.
\item Turn 3: action=finalize, phase=physics, complete=True.
\item Missing fields: none.
\item LLM used=True; fallback=none; backend=llm\_semantic\_frame+runtime\_semantic+confirmed\_pending\_overwrite.
\end{itemize}
\subsection{Case MTM16 (nested\_target)}
\textbf{Dialogue:}
\begin{itemize}
\item User: 我想在一个盒子里嵌一个圆柱靶，先别急着补全。
\item Assistant: 好的，我们已确认了您想要的几何结构类型（盒中嵌圆柱）。接下来还需要您提供以下信息，以便继续配置：

1.  **几何**：外层盒子的具体尺寸，以及内嵌圆柱的半径。
2.  **材料**：为盒体和圆柱体分别指定材料。
3.  **源**：设置粒子源的类型、粒子种类和能量。
4.  **物理**：选择模拟所需的物理过程列表。
5.  **输出**：设定输出数据的格式和保存路径。

请随时告诉我其中任何一项的具体内容。
\item User: 外盒 120 mm x 120 mm x 120 mm，内嵌铜圆柱半径 20 mm、半长 40 mm，嵌套间隙 2 mm；材料 ?；束流 gamma 1 MeV，位置 (0,0,-150) mm，方向 +z；物理列表 QBBC；输出 json。
\item Assistant: 我注意到您新提供的几何参数会覆盖之前已确认的设置。具体来说，外层盒子的尺寸、内嵌圆柱的半径和半长，以及嵌套间隙都将被更新。

如果您确定要应用这些新数值，请直接回复“确认”；如果您希望保留原来的参数，请回复“保持原值”。
\item User: 确认
\item Assistant: 好的，配置已完成。目前我们已确认的参数包括：嵌套间隙、内嵌圆柱半长、内嵌圆柱半径以及外层容器的尺寸。
\end{itemize}
\textbf{Internal Chain (Natural Language):}
\begin{itemize}
\item Turn 1: action=summarize\_progress, phase=materials, complete=False.
\item Missing fields: materials.volume\_material\_map, source.type, source.particle, source.energy, source.position, source.direction, physics.physics\_list, output.format, output.path, geometry.ask.nest.parent\_size, geometry.ask.nest.child\_radius, geometry.ask.nest.child\_half\_length, geometry.ask.nest.clearance.
\item LLM used=True; fallback=none; backend=llm\_slot\_frame+runtime\_semantic.
\item Turn 2: action=confirm\_overwrite, phase=source, complete=False.
\item Missing fields: none.
\item LLM used=True; fallback=none; backend=llm\_slot\_frame+runtime\_semantic.
\item Turn 3: action=finalize, phase=physics, complete=True.
\item Missing fields: none.
\item LLM used=True; fallback=none; backend=llm\_semantic\_frame+runtime\_semantic+confirmed\_pending\_overwrite.
\end{itemize}
\subsection{Case MTM17 (coaxial\_shield)}
\textbf{Dialogue:}
\begin{itemize}
\item User: I need a shell-style shielding toy model, but the shell thicknesses are still open.
\item Assistant: Geometry: updated geometry structure; still needs shell inner radius, shell thicknesses. Materials: still needs volume-to-material binding. Source: updated source energy, source position, source direction; still needs particle type. Physics: updated physics list. Output: still needs output format, output path.
\item User: Use a shell with inner radius 15 mm, thicknesses 5 mm and 8 mm, half length 40 mm, nested core radius 10 mm, nested core half length 30 mm, clearance 1 mm, material lead; isotropic neutron source 2 MeV at (0,0,0) mm; physics Shielding; output hdf5.
\item Assistant: This would overwrite an already confirmed value. The pending changes are:
- Source type: from "point" to "isotropic"
- Source energy: from 1.0 MeV to 2.0 MeV

Please reply with **'confirm'** to apply these changes, or **'keep existing'** to keep the current values.
\item User: confirm
\item Assistant: Configuration complete. The following items are confirmed: source position, source energy, source type, and shell clearance.
\end{itemize}
\textbf{Internal Chain (Natural Language):}
\begin{itemize}
\item Turn 1: action=summarize\_progress, phase=materials, complete=False.
\item Missing fields: materials.volume\_material\_map, source.particle, output.format, output.path, geometry.ask.shell.inner\_radius, geometry.ask.shell.thicknesses, geometry.ask.shell.half\_length, geometry.ask.shell.child\_radius, geometry.ask.shell.child\_half\_length, geometry.ask.shell.clearance.
\item LLM used=True; fallback=none; backend=llm\_slot\_frame+runtime\_semantic.
\item Turn 2: action=confirm\_overwrite, phase=source, complete=False.
\item Missing fields: none.
\item LLM used=True; fallback=none; backend=llm\_slot\_frame+runtime\_semantic.
\item Turn 3: action=finalize, phase=physics, complete=True.
\item Missing fields: none.
\item LLM used=True; fallback=none; backend=llm\_semantic\_frame+runtime\_semantic+confirmed\_pending\_overwrite.
\end{itemize}
\subsection{Case MTM18 (boolean\_cavity)}
\textbf{Dialogue:}
\begin{itemize}
\item User: 我想做一个挖空的盒体几何，但先只给你意图。
\item Assistant: 好的，目前进度如下：

**已设置**：几何结构类型（挖空盒体）。
**待补充**：为了完成挖空操作，还需要您提供两个部件的具体尺寸（部件A的x、y、z尺寸，以及部件B的x、y尺寸）。

此外，材料、粒子源、物理过程和输出设置等信息也仍待补充。我们一步步来。
\item User: 做减法布尔体：A 为 40 mm x 40 mm x 40 mm 盒体，B 为 20 mm x 20 mm x 20 mm 盒体，从 A 中减去 B；材料 铝；点源 gamma 0.3 MeV，位置 (0,0,-60) mm，方向 +z；物理列表 QBBC；输出 xml。
\item Assistant: 我检测到这会覆盖已经确认的值。待修改内容：geometry.graph\_program：\{'nodes': [\{'id': 'a', 'type': 'Box', 'x': 29.287457362487178, 'y': 16.313688044337646, 'z': 53.82008547798903\}, \{'id': 'b', 'type': 'Box', 'x': 10.432721500065707, 'y': 45.19115032300169, 'z': 32.42666876844392\}, \{'id': 'bool', 'type': 'BooleanBinary', 'op': 'subtraction', 'left': 'a', 'right': 'b'\}], 'root': 'boolean', 'constraints': []\} -> \{'nodes': [\{'id': 'a', 'type': 'Box', 'x': 40.0, 'y': 40.0, 'z': 40.0\}, \{'id': 'b', 'type': 'Box', 'x': 20.0, 'y': 20.0, 'z': 20.0\}, \{'id': 'bool', 'type': 'BooleanBinary', 'op': 'subtraction', 'left': 'a', 'right': 'b'\}], 'root': 'boolean', 'constraints': []\}。如果要应用这些修改，请直接回复“确认”；如果要保留原值，请直接说“保持原值”。
\item User: 确认
\item Assistant: 配置已完成。当前已确认的设置有：布尔 solid A 和 B 的尺寸、输出路径以及输出格式。
\end{itemize}
\textbf{Internal Chain (Natural Language):}
\begin{itemize}
\item Turn 1: action=summarize\_progress, phase=geometry, complete=False.
\item Missing fields: geometry.params.bool\_a\_x, geometry.params.bool\_a\_y, geometry.params.bool\_a\_z, geometry.params.bool\_b\_x, geometry.params.bool\_b\_y, geometry.params.bool\_b\_z, materials.volume\_material\_map, source.type, source.particle, source.energy, source.position, source.direction, physics.physics\_list, output.format, output.path, geometry.ask.boolean.solid\_a\_size, geometry.ask.boolean.solid\_b\_size.
\item LLM used=True; fallback=none; backend=llm\_slot\_frame+runtime\_semantic.
\item Turn 2: action=confirm\_overwrite, phase=materials, complete=False.
\item Missing fields: none.
\item LLM used=True; fallback=none; backend=llm\_slot\_frame+runtime\_semantic.
\item Turn 3: action=finalize, phase=source, complete=True.
\item Missing fields: none.
\item LLM used=True; fallback=none; backend=llm\_semantic\_frame+runtime\_semantic+confirmed\_pending\_overwrite.
\end{itemize}
\subsection{Case MTM19 (ring\_complete)}
\textbf{Dialogue:}
\begin{itemize}
\item User: Create a complete ring setup: 16 modules, each 6 mm x 6 mm x 2 mm, radius 32 mm, clearance 0.5 mm, material silicon; point gamma source 0.511 MeV at (0,0,0) mm; physics QBBC; output json.
\item Assistant: Your configuration is complete. I have set up the ring geometry with 16 silicon modules, a 0.511 MeV gamma point source at the origin, and the QBBC physics list. Output will be in JSON format.

The following items are confirmed: volume-to-material binding, output path, clearance, and radius.
\end{itemize}
\textbf{Internal Chain (Natural Language):}
\begin{itemize}
\item Turn 1: action=finalize, phase=materials, complete=True.
\item Missing fields: none.
\item LLM used=True; fallback=none; backend=llm\_slot\_frame+runtime\_semantic.
\end{itemize}
\subsection{Case MTM20 (grid\_complete)}
\textbf{Dialogue:}
\begin{itemize}
\item User: 请创建完整的阵列配置：3 x 3 盒体阵列，每个模块 12 mm x 12 mm x 3 mm，pitch\_x 15 mm，pitch\_y 15 mm，间隙 1 mm，材料 碘化铯；点源 gamma 0.662 MeV，位置 (0,0,-90) mm，方向 +z；物理列表 QBBC；输出 root。
\item Assistant: 配置已完成。当前已确认的设置包括：体积与材料绑定、输出路径、阵列间隙（clearance）、以及 y 方向和 x 方向的阵列间距（pitch）。
\end{itemize}
\textbf{Internal Chain (Natural Language):}
\begin{itemize}
\item Turn 1: action=finalize, phase=materials, complete=True.
\item Missing fields: none.
\item LLM used=True; fallback=none; backend=llm\_slot\_frame+runtime\_semantic.
\end{itemize}
\subsection{Case MTM21 (layered\_detector)}
\textbf{Dialogue:}
\begin{itemize}
\item User: I want a stacked detector slab, but the individual layer thicknesses are not fixed yet.
\item Assistant: Here's what we have so far:

**Geometry:** You've set up a stacked detector structure. We still need the footprint (size) of each layer and their individual thicknesses.

**Materials:** We need to assign materials to the detector volumes.

**Source:** We need the source type, particle type, and its energy and position.

**Physics \& Output:** We still need to select the physics list and specify the output format and file path.
\item User: Use stack footprint 25 mm x 25 mm, thicknesses 1 mm, 2 mm, 1 mm, layer clearance 0.2 mm, outer box 30 mm x 30 mm x 12 mm, container clearance 0.5 mm, material silicon; beam electron 5 MeV from (0,0,-50) mm along +z; physics FTFP\_BERT; output json.
\item Assistant: Your new geometry settings would overwrite previously confirmed values.

**Pending Change:**
- **Stack footprint:** 25 mm x 25 mm (was \textasciitilde{}16.3 mm x \textasciitilde{}10.3 mm)
- **Stack layer thicknesses:** 1 mm, 2 mm, 1 mm (was \textasciitilde{}6.68 mm, \textasciitilde{}1.19 mm, \textasciitilde{}5.59 mm)
- **Layer clearance:** 0.2 mm (was \textasciitilde{}0.73 mm)
- **Outer container size:** 30 mm x 30 mm x 12 mm (was \textasciitilde{}18.1 mm x \textasciitilde{}45.5 mm x \textasciitilde{}16.7 mm)
- **Container clearance:** 0.5 mm (was \textasciitilde{}0.87 mm)

Reply **'confirm'** to apply these new values, or **'keep existing'** to keep the current geometry.
\item User: confirm
\item Assistant: Your configuration is complete. The following geometry items are confirmed: stack container clearance, stack container size, stack layer clearance, and stack layer thicknesses.
\end{itemize}
\textbf{Internal Chain (Natural Language):}
\begin{itemize}
\item Turn 1: action=summarize\_progress, phase=materials, complete=False.
\item Missing fields: materials.volume\_material\_map, source.type, source.particle, source.energy, source.position, source.direction, physics.physics\_list, output.format, output.path, geometry.ask.stack.footprint, geometry.ask.stack.thicknesses, geometry.ask.stack.layer\_clearance, geometry.ask.stack.parent\_size, geometry.ask.stack.nest\_clearance.
\item LLM used=True; fallback=none; backend=llm\_semantic\_frame+runtime\_semantic.
\item Turn 2: action=confirm\_overwrite, phase=source, complete=False.
\item Missing fields: none.
\item LLM used=True; fallback=none; backend=llm\_slot\_frame+runtime\_semantic.
\item Turn 3: action=finalize, phase=physics, complete=True.
\item Missing fields: none.
\item LLM used=True; fallback=none; backend=llm\_semantic\_frame+runtime\_semantic+confirmed\_pending\_overwrite.
\end{itemize}
\subsection{Case MTM22 (embedded\_sensor)}
\textbf{Dialogue:}
\begin{itemize}
\item User: 请创建完整的嵌套配置：外盒 80 mm x 80 mm x 80 mm，内嵌铅圆柱半径 15 mm、半长 25 mm，嵌套间隙 1 mm；材料 ?；点源 gamma 1 MeV，位置 (0,0,-100) mm，方向 +z；物理列表 Shielding；输出 csv。
\item Assistant: 配置已完成。当前已确认的设置包括：体积与材料绑定、输出路径、嵌套间隙（clearance）以及圆柱半径。
\end{itemize}
\textbf{Internal Chain (Natural Language):}
\begin{itemize}
\item Turn 1: action=finalize, phase=materials, complete=True.
\item Missing fields: none.
\item LLM used=True; fallback=none; backend=llm\_slot\_frame+runtime\_semantic.
\end{itemize}
\subsection{Case MTM23 (layered\_shield\_complete)}
\textbf{Dialogue:}
\begin{itemize}
\item User: Create a complete shell configuration: inner radius 12 mm, thicknesses 4 mm and 6 mm, half length 35 mm, nested core radius 8 mm, nested core half length 25 mm, clearance 1 mm, material iron; isotropic gamma source 0.8 MeV at (0,0,0) mm; physics QBBC; output hdf5.
\item Assistant: Your configuration is complete. The following items are confirmed: box dimensions (X, Y, Z) and volume-to-material bindings.
\end{itemize}
\textbf{Internal Chain (Natural Language):}
\begin{itemize}
\item Turn 1: action=finalize, phase=materials, complete=True.
\item Missing fields: none.
\item LLM used=True; fallback=none; backend=llm\_slot\_frame+runtime\_semantic.
\end{itemize}
\subsection{Case MTM24 (boolean\_complete)}
\textbf{Dialogue:}
\begin{itemize}
\item User: 请创建完整布尔配置：做并集布尔体，A 为 30 mm x 30 mm x 20 mm 盒体，B 为 10 mm x 10 mm x 40 mm 盒体，与 A 做 union；材料 铝；束流 gamma 0.5 MeV，位置 (0,0,-70) mm，方向 +z；物理列表 QBBC；输出 json。
\item Assistant: 配置已完成。当前已确认的设置有：盒体在 x、y、z 方向的尺寸，以及体积与材料的绑定。
\end{itemize}
\textbf{Internal Chain (Natural Language):}
\begin{itemize}
\item Turn 1: action=finalize, phase=materials, complete=True.
\item Missing fields: none.
\item LLM used=True; fallback=none; backend=llm\_slot\_frame+runtime\_semantic.
\end{itemize}
\end{document}